\section{Conclusion}

This systematic review examined thirteen studies on deep learning approaches for ambulatory seizure monitoring using non-EEG biosignals, encompassing both detection and forecasting applications published between 2013 and 2026. The evidence reveals a field at a crossroads: detection technology has achieved substantial technical maturity for convulsive seizures, yet critical barriers prevent widespread clinical adoption. Below, we synthesize the current state of methodological development, identify persistent translation challenges, and outline priorities for future research.

\subsection{Current State of the Field}

The field of ambulatory seizure monitoring using non-EEG biosignals has achieved substantial technical maturity in detection capabilities while forecasting remains an emerging challenge. Across nine detection studies, deep learning systems demonstrate high sensitivity ranging from 71.6\% to 100\% for identifying convulsive seizures, with accelerometry-based approaches proving particularly effective for motor manifestations. This performance foundation establishes that wearable devices can reliably detect generalized tonic-clonic and other convulsive seizure types in controlled settings. However, false alarm rates reveal a critical translation gap: only two studies approach the ILAE Phase 3 benchmark for validated devices (see Section~\ref{sec:ilae-standards}). Fine et al. report 0.023/h for tonic seizures, while Dong et al. achieve 0.165/h for nocturnal severe seizures, representing the only systems approaching deployable thresholds.

Seizure forecasting exhibits markedly less maturity than detection. Across four prospective studies, only 43.5\% to 83\% of patients demonstrate better-than-chance performance, highlighting a fundamental responder problem that limits universal applicability. This patient heterogeneity means that aggregated performance metrics obscure substantial individual variation. When rigorous leave-one-subject-out validation is applied, as in Meisel et al., fewer than half of patients achieve significant forecasting capability. The forecasting challenge thus appears less tractable than detection in the current state of development.

Sensor modality effectiveness follows a clear hierarchy. Accelerometry dominates the evidence base, appearing in eight of thirteen studies and providing the most robust signal for convulsive seizure detection. Cardiac signals contribute value in six studies, particularly for preictal changes and forecasting applications, but standalone electrocardiogram approaches suffer from excessive motion artifact sensitivity and false alarm rates. Multimodal fusion strategies, combining accelerometry with cardiac or electrodermal measures in four recent studies, represent the current frontier for improving specificity while maintaining sensitivity.

The clinical utility of current systems remains tightly bounded by seizure semiology. Convulsive seizures with prominent motor components are highly detectable through wrist-worn accelerometry, achieving sensitivity above 90\% in multiple studies. By contrast, non-motor seizures---including focal aware seizures and absence seizures---remain largely undetectable through available wearable modalities. This detection gap creates a fundamental limitation: patients whose seizure semiology lacks convulsive features cannot benefit from current ambulatory monitoring approaches, regardless of algorithmic sophistication.

Validation rigor represents another limiting factor in the field's current state. Only four of thirteen studies employ leave-one-subject-out or patient-independent cross-validation, with most relying on retrospective analyses that likely overestimate real-world performance. Prospective validation remains rare, and multi-center evidence is limited. The evidence base thus reflects promising but early-stage technology, where algorithmic innovation has outpaced rigorous clinical validation in diverse ambulatory settings.

\subsection{Methodological Evolution}

The methodological landscape of wearable seizure detection and forecasting has undergone substantial transformation from 2013 to 2026, reflecting both technical advances in deep learning and increasing rigor in clinical validation. This evolution occurred in three interconnected dimensions: algorithmic sophistication, validation methodology, and clinical readiness assessment.

Early work from 2013 to 2015 relied predominantly on traditional machine learning approaches. Studies during this period employed k-nearest neighbors, artificial neural networks, and support vector machines operating on handcrafted features extracted from cardiac signals. These methods, while interpretable, required extensive domain expertise for feature engineering and typically achieved moderate performance on small datasets.

The period from 2016 to 2019 marked a transition toward heart rate variability-based deep learning models, particularly multilayer perceptrons and long short-term memory networks. This shift enabled automated feature extraction from temporal HRV patterns, though these studies often relied on patient-specific hyperparameter tuning that inflated apparent performance. The architectural advances during this era prioritized sensitivity over specificity, resulting in unacceptably high false alarm rates.

From 2020 to 2022, algorithm development emphasized architectural complexity through hybrid models combining LSTM networks with one-dimensional convolutional layers, alongside ensemble methods. These architectures explicitly addressed temporal dependencies in cardiac and accelerometric signals while improving robustness through ensemble diversity. The introduction of leave-one-subject-out cross-validation during this period revealed that patient-specific tuning had previously obscured substantial performance variation across individuals.

The most recent period from 2023 to 2026 has witnessed the emergence of attention-based architectures, matrix profile anomaly detection, and self-attentive autoencoders. Attention mechanisms enable models to selectively weight informative signal segments while suppressing motion artifacts and noise. Unsupervised approaches have demonstrated that seizure-related autonomic changes can be detected without requiring large annotated seizure datasets, potentially addressing the chronic data scarcity problem.

Validation methodology has evolved more slowly than algorithm architecture, creating a persistent gap between technical performance and clinical applicability. Early studies relied on simple train-test splits without subject-level separation, enabling data leakage that severely overestimated real-world performance. The recognition of this problem led to widespread adoption of LOSO cross-validation, which provides a more realistic assessment of generalizability to new patients. Despite this improvement, prospective validation remains rare, with most studies relying on retrospective analyses.

The International League Against Epilepsy Phase framework provides a structured approach to assessing clinical readiness. The current evidence base remains concentrated in Phases 1 and 2, with only two studies achieving Phase 2 validation. No studies have progressed to Phase 3 or beyond, indicating that wearable seizure detection and forecasting remain in early translational stages despite more than a decade of research.

A persistent obstacle to progress is the extreme heterogeneity in performance metrics reporting. Studies report sensitivity, specificity, accuracy, precision, recall, F1-score, and area under the ROC curve with no consensus on primary endpoints. False alarm rates, the most clinically relevant metric for patient adherence, are reported in incompatible units or omitted entirely. This diversity impedes direct comparison across studies and meta-analytic synthesis. The absence of standardized benchmark datasets further compounds this problem.

The relationship between algorithmic sophistication and clinical readiness follows a counterintuitive trajectory. Architectural advances have consistently outpaced validation rigor, with studies employing state-of-the-art deep learning models yet failing to demonstrate prospective validity. This disconnect suggests that research effort has been misallocated toward incremental architectural improvements rather than addressing fundamental validation challenges.

\subsection{Barriers to Clinical Translation}

Four interrelated barriers prevent widespread clinical adoption of deep learning systems for ambulatory seizure monitoring: false alarm rates, methodological limitations in validation studies, patient heterogeneity, and scope restrictions regarding detectable seizure types.

The most immediate technical barrier is the unacceptable false alarm rate reported in most studies. While detection sensitivity consistently exceeds 90\% for convulsive seizures, only two of thirteen detection studies approach the ILAE Phase 3 benchmark for validated devices (see Section~\ref{sec:ilae-standards}). The remaining systems generate 20 to 200 times more false alarms than clinicians and patients can accept, with ECG-based approaches reaching 1.91 to 39.75 false alarms per hour. This discrepancy means that most validated algorithms would generate 50 to 500 false alarms per day in home settings, ensuring patient non-adoption and caregiver fatigue. The single exception among cardiac-based systems is Fine's tonic seizure detector, which demonstrates that aggressive artifact rejection and narrow seizure focus can achieve clinical viability, albeit at the cost of limited generalizability.

Three systematic methodological weaknesses inflate reported performance and limit generalizability. First, sample sizes are critically underpowered, with median cohort sizes of 15 to 70 patients. These small samples preclude meaningful subgroup analysis by seizure type, medication regime, or comorbidities. Second, monitoring duration is typically limited to days or weeks rather than months, which biases performance estimates toward common seizure types while underestimating challenges with rare events. Third, setting bias is pervasive: eleven of thirteen studies conduct validation in epilepsy monitoring units under controlled conditions that minimize motion artifacts, environmental noise, and device compliance problems. The two studies that tested systems in actual home environments both reported performance degradation compared to EMU results, suggesting that current validation standards substantially overestimate real-world utility.

The most fundamental barrier to clinical translation is the responder problem: current forecasting and detection approaches work well for a minority of patients but fail for the majority. In leave-one-subject-out cross-validation, which represents true generalization to unseen patients, only 43.5\% of patients achieved better-than-chance forecasting performance. Similar patterns appear in detection studies, where patient success rates range from 44\% to 100\% across studies. This heterogeneity means that aggregated performance metrics obscure the reality that clinicians cannot predict in advance which patients will benefit from wearable monitoring. The absence of baseline biomarkers to identify likely responders forces a trial-and-error approach that is incompatible with clinical practice.

A final barrier is the restricted seizure type coverage of current non-EEG approaches. Accelerometry-based systems reliably detect convulsive and motor seizures but cannot identify focal aware seizures, absence seizures, or other non-motor events that constitute the majority of seizures in many patients. Cardiac-based approaches show promise for preictal detection but are limited by motion artifact sensitivity and the fact that only 60 to 70\% of patients exhibit reliable heart rate changes before seizures. This scope limitation means that even a perfectly optimized wearable system would remain incomplete without complementary EEG or multimodal integration, restricting clinical utility to a subset of the epilepsy population.

\subsection{Future Directions}

The transition from research prototypes to clinically viable seizure monitoring systems requires addressing several methodological and implementation gaps. Based on our findings, we identify three priority areas that warrant immediate attention.

Standardization and benchmarking represent the foundation for progress. The field lacks consensus on evaluation metrics and reporting standards, making cross-study comparisons difficult. Minimum reporting requirements should include patient-level confusion matrices, exact seizure durations, latency distributions, and battery consumption metrics. Open benchmark datasets with annotated non-motor seizures would enable meaningful algorithm comparisons. Additionally, standardized protocols for simulating real-world conditions---sensor displacement, signal artifacts, and data loss---would better prepare systems for deployment.

Real-world deployment challenges must be addressed prospectively. Most validated algorithms remain unevaluated in home environments. Prospective studies must assess long-term reliability across diverse seizure types, daily activities, and patient populations. Edge optimization represents a critical bottleneck: while deep learning models demonstrate strong performance, computational demands often exceed embedded device constraints. Research should explore model compression techniques, lightweight architectures, and energy-efficient inference strategies. Integration with wearable hardware requires addressing signal synchronization, multi-modal data fusion, and continuous operation over extended periods.

Personalization and interpretability are essential for clinical acceptance. Universal models rarely achieve optimal performance across all patients due to high inter-individual variability. Federated learning approaches offer a pathway to personalized models without centralizing sensitive data, preserving privacy while enabling adaptation. Simultaneously, explainable AI methods must bridge the gap between high accuracy and clinical trust. Visual explanations of feature contributions, uncertainty quantification, and interpretable error analysis would facilitate clinician acceptance and regulatory approval.

Expanding detection scope beyond convulsive seizures remains a critical unmet need. Current research predominantly addresses motor seizures, leaving absence seizures and focal impaired awareness seizures largely unaddressed. Future work should develop multi-modal approaches that combine physiological signals with behavioral markers. Hybrid systems integrating wearable data with patient-specific information---medication schedules, seizure diaries, and comorbidities---may improve sensitivity while reducing false alarms.

The path forward requires collaboration between machine learning researchers, clinicians, and industry partners. Priority should be given to prospective validation studies that report both technical and patient-centered outcomes, ensuring that advances translate into meaningful improvements in epilepsy management.
